% answers_short.tex —— 2008 年·全国大纲卷Ⅰ·理科 · 简短答案速查
% 由 tools/make_short_answers.py 自 answers.tex 的 \gksolhead/\gkanswer 提取，
% 并与 2008/2008_数学真题_全国1理_答案要点.md 逐题比对后生成，勿手改（改源 MD 后重跑）。

\begin{gknotice}
\textbf{使用说明}：本节只给最终答案便于快速核对，完整推导见解析版。答案为编写者独立推导，\textbf{非官方标准答案}。
\end{gknotice}

\gkpart{选择题}
\begin{center}\begin{tabular}{ccccccccccccc}
\toprule
题号 & 1 & 2 & 3 & 4 & 5 & 6 & 7 & 8 & 9 & 10 & 11 & 12 \\
\midrule
答案 & C & A & A & D & C & B & A & A & D & D & B & B \\
\bottomrule
\end{tabular}\end{center}

\gkpart{填空题}
\begin{itemize}[leftmargin=2.6em]
  \item[\textbf{13.}] $9$
  \item[\textbf{14.}] $2$
  \item[\textbf{15.}] $\dfrac38$
  \item[\textbf{16.}] $\dfrac16$
\end{itemize}

\gkpart{解答题}
\begin{itemize}[leftmargin=2.6em]
  \item[\textbf{17.}] (1) $\tan A\cot B=4$；(2) $\tan(A-B)$ 的最大值为 $\dfrac34$
  \item[\textbf{18.}] (1) $p=0.001$；(2) 每位投保人的最低保费为 $15$ 元
  \item[\textbf{19.}] (1) 见解析版；(2) $\cos\theta=\dfrac1{\sqrt{10}}$
  \item[\textbf{20.}] (1) $a_n$ 的通项由 $b_n=(a-3)2^{n-1}$ 给出；(2) $a\geqslant-9$
  \item[\textbf{21.}] (1) $e=\dfrac{\sqrt5}{2}$；(2) 双曲线方程为 $\dfrac{x^2}{36}-\dfrac{y^2}{9}=1$
  \item[\textbf{22.}] (1) $f(x)$ 在 $\cos x>-\dfrac12$ 的区间上递增，在 $\cos x<-\dfrac12$ 的区间上递减；(2) $a\geqslant\dfrac13$
\end{itemize}
